\section{Methodology}
\label{sec:methodology}

Our methodology builds an automatic schematic information extraction pipeline for circuit images. The pipeline takes a schematic image as input and extracts structured circuit artifacts, including component locations, component classes, text regions, wire segments, component orientations, electrical nodes, node-component contacts, incidence information, and a SPICE-style netlist when the extracted connectivity is complete. Although netlist generation is the most direct downstream output in the current implementation, the pipeline is designed as a general information extraction system whose intermediate artifacts can also support future extraction of component labels, parameter values, named ports, and other schematic metadata.

The front half of the overall project contains two parts. The first part is the automatic information extraction pipeline. This pipeline can run independently and produce schematic-to-structure outputs without manual intervention. The second part is an optional manual annotation and correction tool. This tool is not required by the automatic pipeline; it is provided to help users inspect intermediate results, fix hard cases, and build cleaner corrected datasets.

\subsection{Automatic Information Extraction Pipeline}
\label{subsec:auto-info-extraction}

The automatic pipeline is organized into six major stages: component detection, text and component masking, line detection, orientation detection, node construction, and incidence/netlist generation. The pipeline is modular by design. Each stage produces explicit intermediate artifacts, so later stages can be inspected, corrected, or rerun without repeating the entire process.

The main workflow is as follows. First, the pipeline detects circuit components and records their classes and bounding boxes. Second, it masks detected components and removes detected text regions to create a cleaner image for wire extraction. Third, it detects wire segments using HAWPv3 and supplements them with pixel-level residual line recovery. Fourth, it predicts orientations for orientation-sensitive components. Fifth, it groups wire segments into electrical nodes using junction and jump detection. Finally, it assigns nodes to component pins, checks the incidence result, and writes a netlist only when the connectivity is valid.

\subsection{Component Detection}
\label{subsec:component-detection}

The component detection stage uses a YOLO11s object detector. The detector outputs component bounding boxes, component classes, and confidence scores. It is trained on a mixed dataset consisting of 3,003 real-world schematic images and 2,000 synthetic schematic images. The real-world data is collected from AMS Net and Image2Net and manually cleaned before training. The synthetic data is generated from Netlistify and is used to improve coverage of component types, orientations, and layout variations.

The trained model detects 18 classes:
\texttt{gnd}, \texttt{vdd}, \texttt{nmos}, \texttt{nmos-bulk}, \texttt{pmos}, \texttt{pmos-bulk}, \texttt{npn}, \texttt{pnp}, \texttt{resistor}, \texttt{capacitor}, \texttt{inductor}, \texttt{diode}, \texttt{voltage\_src}, \texttt{ac\_src}, \texttt{current\_src}, \texttt{battery}, \texttt{amplifier}, and \texttt{switch\_ideal}. These classes are chosen to preserve circuit semantics needed by later stages. For example, separating NMOS and PMOS, and distinguishing MOS devices with explicit bulk terminals, makes the final pin assignment more precise than a generic transistor label.

The detected component boxes are used in two ways. First, they define the component set for the later incidence representation. Second, they are used to mask component interiors before line detection. This masking is important because internal symbol strokes can otherwise be misclassified as wires.

\subsection{Text and Component Masking}
\label{subsec:masking}

Before line detection, the pipeline creates a component-masked image by erasing the detected component regions. It also applies OCR to locate text annotations, but the text is not removed from the image used by HAWPv3. Instead, the OCR output is recorded as text bounding boxes. This keeps the original visual context available for line detection while giving later stages explicit information about where labels, parameter values, and node names appear.

The recorded text boxes are used in two later places. First, during residual black-pixel line enhancement, text regions are masked out so that strokes from labels are not added as extra wire segments. Second, during incidence construction, the text boxes provide cross-validation information. If a candidate touch point or extra line lies inside or very close to a text region, the pipeline can treat it as suspicious and remove it from the connectivity evidence. This helps prevent labels from becoming false wires or false pins without suppressing text before HAWPv3 inference.

\subsection{Line Detection}
\label{subsec:line-detection}

The line detection stage uses HAWPv3, a Holistically-Attracted Wireframe Parsing model, as the main wire detector. HAWPv3 was originally designed for wireframe parsing in natural images, but it works well for schematic line detection because schematic wires are dominated by straight structural segments. We use the pretrained HAWPv3 model directly, so this stage does not require task-specific wire annotation or additional training.

We select HAWPv3 because the node construction stage benefits more from high line recall than from overly conservative line prediction. If the detector finds extra line segments, the pipeline can remove or ignore many of them later using component masks, OCR masks, geometric filtering, and node construction rules. However, if a critical wire segment is never detected, it is much harder to recover the missing connection downstream. Therefore, we prefer a line detector that identifies as many plausible wire segments as possible.

The raw HAWPv3 output is filtered and refined for schematic use. Candidate lines are clipped against detected component boxes so that internal component strokes are removed while external wire stubs are preserved. Lines that are too short, lie mostly in whitespace, or are too close to OCR-detected text are discarded. The pipeline then performs residual black-pixel recovery outside component and text regions. Specifically, it binarizes the image, keeps only black pixels outside the component and OCR text boxes, applies a small morphological closing operation to connect tiny gaps, skeletonizes the remaining strokes, and runs a probabilistic Hough transform to convert the residual pixels into line segments. This residual pass captures short or faint wire fragments that HAWPv3 may miss. The final line set is the union of the filtered HAWPv3 lines and the valid residual pixel-derived lines.

\subsection{Orientation Detection}
\label{subsec:orientation-detection}

The orientation detection stage is used for component classes whose pin semantics depend on direction. This includes MOSFET classes, BJT classes, diodes, and amplifiers. For each detected component in these classes, the pipeline crops the component from the original image and classifies its orientation as up, right, down, or left.

This stage uses a ResNet18 classifier trained on a mixture of real and synthetic component crops. The real crops are collected from Image2Net, while the synthetic crops provide controlled rotations and balanced orientation coverage. The output orientation is saved as a separate artifact and is later used by the incidence construction stage. For example, the orientation determines which side of a MOSFET corresponds to the gate and how the remaining contacts should be assigned to source, drain, or bulk.

\subsection{Junction and Jump Detection}
\label{subsec:junction-jump-detection}

Line detection alone does not determine electrical connectivity. Two wires may visually cross without being connected, while a junction marker indicates that intersecting wires should belong to the same node. To handle this, the pipeline uses a second YOLO11s model to detect junctions and jumps in the component-masked image.

This junction/jump detector is trained on 5,000 synthetic images. Synthetic data is appropriate here because junctions and jumps are local geometric symbols that can be generated with accurate labels and varied surrounding wire patterns. The detector outputs two types of regions. A junction region means that all intersecting lines in that region should be merged into the same electrical node. A jump region means that crossing lines should be assigned to different nodes, while line fragments along the same axis should still be merged as the same visual wire.

The pipeline also handles implicit jumps. If a schematic style uses junction markers, then crossings without a junction marker are treated as non-contacting crossings. In this case, crossing lines at different angles are assigned to different nodes. This rule allows the pipeline to distinguish real electrical connections from visual crossings even when no explicit jump arc is drawn.

\subsection{Node Construction}
\label{subsec:node-construction}

Node construction groups recovered wire segments into electrical nodes. The pipeline first finds line segments whose closest distance is below a node grouping threshold. These pairs are candidates for merging, but the final decision depends on junction and jump evidence.

\begin{figure}[t]
\centering
\fbox{%
\begin{minipage}{0.94\linewidth}
\textbf{Algorithm 1: Junction-aware node construction}

\textbf{Input:} wire segments $\mathcal{L}$, junction regions $\mathcal{J}$, jump regions $\mathcal{U}$, grouping threshold $\tau$.

\textbf{Output:} node partition $\mathcal{V}$ and assignment $\nu:\mathcal{L}\rightarrow\mathcal{V}$.

\begin{enumerate}
    \item Build the close-pair set
    \[
        \mathcal{E}_{c}=\{(\ell_i,\ell_j)\mid d(\ell_i,\ell_j)<\tau\}.
    \]
    \item Initialize union-find over $\mathcal{L}$ and force-merge all lines intersecting the same junction region:
    \[
        \forall J\in\mathcal{J},\quad
        \mathrm{Union}\bigl(\{\ell\in\mathcal{L}\mid \ell\cap J\neq\emptyset\}\bigr).
    \]
    \item For each $(\ell_i,\ell_j)\in\mathcal{E}_{c}$, estimate contact point $p_{ij}$ and angle $\theta_{ij}$.
    \item Build the split set
    \[
        \mathcal{S}=\{(\ell_i,\ell_j)\in\mathcal{E}_{c}\mid
        p_{ij}\in\mathcal{U}\ \text{or}\ \mathrm{ImplicitJump}(p_{ij},\theta_{ij},\mathcal{J})\}.
    \]
    \item Merge all close pairs not marked as split:
    \[
        \forall(\ell_i,\ell_j)\in\mathcal{E}_{c}\setminus\mathcal{S},\quad
        \mathrm{Union}(\ell_i,\ell_j).
    \]
    \item For each explicit jump region $U\in\mathcal{U}$, reconnect same-axis fragments:
    \[
        \mathrm{Union}\bigl(\{\ell\in\mathcal{L}\mid \ell\cap U\neq\emptyset,\ \mathrm{axis}(\ell)=a\}\bigr)
        \quad a\in\{x,y\}.
    \]
    \item Return the connected components of union-find as $\mathcal{V}$ and assign each line $\ell$ to its component $\nu(\ell)$.
\end{enumerate}
\end{minipage}}
\caption{Junction-aware node construction.}
\label{alg:node-construction}
\end{figure}

Algorithm~\ref{alg:node-construction} summarizes the node construction rule. The result is a node ID for each recovered wire segment. This stage is the bridge between visual wire detection and circuit connectivity: it converts many local line segments into a smaller set of electrical nets.

\subsection{Incidence Construction}
\label{subsec:incidence-construction}

After node construction, the pipeline determines how electrical nodes touch components. A touch is recorded when a node endpoint intersects or approaches a component boundary. Since small gaps can be introduced by masking or imperfect line detection, the pipeline also extends endpoints by a bounded distance to recover near-boundary contacts. This step helps recover real connections that are visually present but slightly broken in the detected line geometry.

The incidence construction stage maps node-component touches to component pins. The mapping uses the component class, component orientation, touched edge, and contact position. For two-terminal devices such as resistors, capacitors, and inductors, the assignment is mostly symmetric. For sources, the assignment determines positive and negative terminals. For MOS and BJT devices, the orientation determines the gate or base side, and the remaining contacts are ordered to assign source/drain or collector/emitter. Supply symbols are converted into canonical \texttt{GND} and \texttt{VDD} nets.

The incidence builder also performs limited corrections when the visual evidence is clear. It can merge duplicate touches from the same node, remove extra touches overlapping OCR text, rescue missing pins from nearby line evidence, and handle close component-edge connections. These operations are conservative and are intended to fix recognition artifacts rather than infer unsupported circuitry.

\begin{figure}[t]
\centering
\fbox{%
\begin{minipage}{0.94\linewidth}
\textbf{Algorithm 2: Incidence construction and pin assignment}

\textbf{Input:} components $\mathcal{C}$, orientations $\mathcal{O}$, touches $\mathcal{T}$, OCR text boxes $\mathcal{X}$, wire segments $\mathcal{L}$.

\textbf{Output:} incidence map $\mathcal{A}$, node names $\mathcal{N}$, flags $\mathcal{F}$, red flags $\mathcal{R}$.

\begin{enumerate}
    \item For each component $c_i\in\mathcal{C}$, collect its incident touches
    \[
        \mathcal{T}_i=\{t\in\mathcal{T}\mid \mathrm{bbox}(t)=i\}.
    \]
    \item Build the port-candidate pool
    \[
        \mathcal{P}=\{t\in\mathcal{T}\mid \mathrm{node}(t)\ \text{touches only one component}\}.
    \]
    \item For each two-terminal component with $|\mathcal{T}_i|=3$, keep the best opposite-edge pair and move the redundant touch to $\mathcal{P}$ when valid.
    \item Canonicalize supply nets:
    \[
        \mathrm{node}(t)=\texttt{GND}\ \text{if }t\text{ touches a \texttt{gnd} symbol},\quad
        \mathrm{node}(t)=\texttt{VDD}\ \text{if }t\text{ touches a \texttt{vdd} symbol}.
    \]
    \item For each non-supply component, compute an initial pin map
    \[
        A_i^{0}=\Phi\bigl(\mathrm{class}(c_i),\,\mathcal{O}_i,\,\mathcal{T}_i\bigr),
    \]
    where $\Phi$ uses orientation-aware rules for MOS/BJT pins, diode polarity, source polarity, symmetric passive pins, and ordered amplifier pins.
    \item Refine the pin map with visual-evidence repairs:
    \[
        A_i=\mathrm{Repair}(A_i^{0},\mathcal{T}_i,\mathcal{P},\mathcal{X},\mathcal{L}),
    \]
    including duplicate-touch merging, OCR-overlap pruning, nearby-line rescue, close-component-edge rescue, and black-pixel short merging.
    \item Validate the required pin set:
    \[
        R_i=\mathrm{Validate}\bigl(A_i,\mathrm{RequiredPins}(\mathrm{class}(c_i))\bigr).
    \]
    \item Emit $\mathcal{A}=\{A_i\}$, $\mathcal{N}$, $\mathcal{F}$, and $\mathcal{R}=\{R_i\mid R_i\neq\emptyset\}$.
\end{enumerate}
\end{minipage}}
\caption{Orientation-aware incidence construction from node-component touches.}
\label{alg:incidence-construction}
\end{figure}

\subsection{Red Flags and Netlist Generation}
\label{subsec:red-flags-output}

Before generating a netlist, the pipeline checks whether each component has a valid pin assignment. If a component has missing orientation, too few or too many touches, ambiguous polarity, unexpected diode-side contacts, or a one-touch connection without a valid port candidate, the component receives a red flag.

The red-flag mechanism prevents the pipeline from outputting unreliable netlists. If an image has no component-level red flags, the incidence representation is converted into a SPICE-style netlist. The netlist writer uses standard pin orders, including \texttt{D G S B} for MOS devices, \texttt{C B E} for BJTs, \texttt{A K} for diodes, and \texttt{+ -} for sources. If red flags remain, the pipeline keeps the incidence JSON and visual overlays for inspection instead of forcing a netlist.

\subsection{Optional Manual Annotation Pipeline}
\label{subsec:manual-annotation}

In addition to the automatic information extraction pipeline, we provide a manual fixing interface for reviewing and correcting extraction results. The interface shows the original image, component detections, masked image, orientation overlay, recovered lines, node grouping, node-component touches, and incidence visualization. Images with red flags are prioritized, but users can inspect any processed schematic.

The tool supports two types of correction. First, users can fix many images by only adjusting the parameters in the pipeline configuration. These parameters include component detection thresholds, OCR settings, line filtering thresholds, node grouping thresholds, junction/jump detection settings, and incidence rescue parameters. After a parameter change, the tool reruns the automatic pipeline from the earliest affected stage instead of rerunning the whole pipeline. This makes parameter tuning a fast way to repair images whose errors come from threshold choices rather than wrong manual annotations.

Second, users can directly edit intermediate artifacts. They can modify component boxes and classes, set orientations, adjust line-node assignments, add or delete touches, and annotate junction or jump regions. Manual junction and jump annotations are stored as additive overrides and are combined with automatic detections during node construction. This allows the manual tool to correct difficult cases without replacing the automatic pipeline.

The manual annotation pipeline is therefore an optional enhancement to the automatic extraction system. It improves usability, supports dataset cleaning, and provides a practical way to handle schematics that are too ambiguous or noisy for fully automatic extraction.
